\section{Introduction}

Portfolio risk depends on how stocks move together, especially during periods of
large losses. Historical simulation and copula models provide established ways
to estimate this risk, but their results depend on the available history and
assumptions about dependence. Neural stochastic differential equations (neural
SDEs) can generate joint return paths. Whether those paths lead to more useful
risk estimates is an open question.

This thesis will fit a multivariate neural SDE to the daily returns of a selected
group of OSEBX stocks. It will generate return paths and assess how well they
reproduce features of observed data, including volatility clustering, dependence
between stocks, and joint losses. Applying portfolio weights to the generated
paths will give a distribution of portfolio losses from which value at risk
(VaR) can be estimated.

The neural SDE will be compared with historical simulation, weighted historical
simulation, and Gaussian- and Student's \emph{t}-copula models using the same
portfolio and evaluation period. The aim is to find out whether the neural SDE
adds value for generating realistic scenarios and estimating portfolio risk,
while acknowledging that rare losses are difficult to assess from limited data.
